\chapter{Amazon Web Services jako środowisko chmurowe}
\label{chap:aws}

Amazon Web Services (AWS) jest największą pod względem udziału w~rynku platformą chmury publicznej, oferującą
kilkaset usług. Od surowych zasobów obliczeniowych, przez zarządzane bazy danych i~platformy kontenerowe, po
wyspecjalizowane usługi analityczne. Rozdział ten charakteryzuje platformę AWS oraz omawia usługi, z~których
zbudowano środowisko uruchomieniowe zrealizowanego systemu, a~kończy się porównaniem Amazon EKS z~konkurencyjnymi
zarządzanymi usługami Kubernetes.

\section{Charakterystyka AWS i jego architektury globalnej}
\label{sec:charakterystyka-aws}

Infrastruktura AWS jest podzielona na \textbf{regiony}. Są to geograficznie odseparowane lokalizacje (np. \code{eu-central-1} we Frankfurcie).
Każdy z~nich jest oddzielną domeną awarii i~osobnym miejscem przechowywania danych. Każdy region składa się
z~co najmniej trzech \textbf{stref dostępności} (\emph{Availability Zones}, AZ). Są to fizycznie rozdzielone 
centra danych o~niezależnym zasilaniu i łączności. Połączone są one ze sobą łączami o~niskich opóźnieniach.
Rozmieszczenie komponentów systemu w~wielu strefach dostępności jest podstawową techniką zapewniania
wysokiej dostępności. Awaria pojedynczej strefy nie przerywa działania systemu \cite{awsdocs}.

Model rozliczeń AWS ma wiele form. W podstawowej wersji opłaty naliczane są za faktyczne zużycie zasobów
(czas działania instancji, pojemność, transfer danych). Bez zobowiązań wstępnych. W tym projekcie wykorzystana
została pula bezpłatnych zasobów (\emph{Free Tier}), udostępniona czasowo dla nowych kont. Z tego powodu
wiele decyzji co do doboru parametrów środowiska została podjęta w zakresie bezpłatnego poziomu.

Istotnym elementem charakterystyki platformy jest \textbf{model odpowiedzialności współdzielonej} (\emph{shared
responsibility model}). AWS odpowiada za bezpieczeństwo \emph{samej chmury}. Czyli infrastruktury fizycznej, hipernadzorcy,
oprogramowania usług zarządzanych. Natomiast klient za bezpieczeństwo \emph{w~chmurze}. Czyli konfigurację sieci
i~uprawnień, ochronę danych oraz poprawność własnych aplikacji \cite{awsdocs}.

\section{Amazon EKS - zarządzana usługa Kubernetes}
\label{sec:amazon-eks}

Amazon Elastic Kubernetes Service (EKS) udostępnia klaster Kubernetes. Płaszczyzna sterowania jest w~całości
zarządzana przez AWS. Użytkownik widzi tylko punkt końcowy API \cite{awsdocs}. Użytkownik odpowiada
za dobór wezłów roboczych oraz uruchomienie obciążenia. EKS wykorzystuje niezmodyfikowaną, certyfikowaną 
dystrybucję Kubernetes, dzięki czemu obciążenia i~manifesty pozostają przenośne pomiędzy środowiskami.

Warstwę obliczeniową klastra EKS można zrealizować na dwa sposoby. Jako samodzielnie zarządzane instancje EC2
\textbf{zarządzane grupy węzłów} (\emph{managed node groups}). W~których AWS automatyzuje provisionowanie,
aktualizacje i~wymianę instancji EC2. AWS Fargate, model bezserwerowy, w~którym pojedyncze pody
uruchamiane są bez jawnie widocznych węzłów. W~zrealizowanym systemie wybrano zarządzane grupy węzłów: w~odróżnieniu
od Fargate zachowują one pełną kontrolę nad typem instancji.

Integralną częścią EKS są \textbf{dodatki} (\emph{add-ons}). Zarządzane przez AWS komponenty systemowe klastra.
Przykładem takiej wtyczki może być sieciowa VPC CNI. Przydziela ona podom adresy IP bezpośrednio z~puli
adresowej sieci VPC. Bez niej liczba podów na węźle jest ograniczona liczbą interfejsów sieciowych instancji.
W tym projekcie użyłem małych instancji co stanowiło bariere przy wdrożeniu wszystkich komponentów.

Tożsamość i~dostęp do klastra EKS realizowane są poprzez integrację z~systemem IAM. Zarówno administratorzy, jak
i~węzły robocze uwierzytelniają się wobec API Kubernetes przy użyciu tożsamości IAM (sekcja~\ref{sec:bezpieczenstwo-kubernetes}).

\section{Usługi powiązane wykorzystane w projekcie}
\label{sec:uslugi-powiazane}

\subsection{Amazon EC2 jako warstwa obliczeniowa}
\label{subsec:ec2}

Amazon Elastic Compute Cloud (EC2) dostarcza maszyny wirtualne (\emph{instancje}). Mają one zróżnicowane zasoby
opisane typami instancji. W zrealizowanym systemie użyte zostały jako węzły EKS. Stanowią one warstwę obliczeniową klastra.
W celu zmieszczenia w darmowym programie użyte zostały instancje \code{t3.small} w~trybie \emph{on-demand}.
Grupa węzłów posiada konfigurowalny przedział wielkości (minimalna, pożądana i~maksymalna liczba instancji). 
Stanowi to infrastrukturalny fundament skalowania klastra.

\subsection{Amazon RDS for PostgreSQL jako warstwa danych}
\label{subsec:rds}

Amazon Relational Database Service (RDS) jest zarządzaną usługą relacyjnych baz danych, przejmującą od użytkownika
instalację, łatanie, kopie zapasowe i~odtwarzanie silnika bazodanowego. W~realizowanym systemie użyto RDS PosgreSQL.
Korzystają z~niej wszystkie 3 mikroserwisy. Instancja umieszczona jest w~podsieci prywatnej. Z~grupą bezpieczeństwa,
która umożliwia dostęp tylko z sieci klastra. Decyzja użycia RDS zamiast bazy uruchamianej w~klastrze, wynika
z~dostarczenia wszystkich potrzebnych mechanizmów dla bazy produkcyjnej przez serwis.

\subsection{Amazon ECR - rejestr obrazów kontenerów}
\label{subsec:ecr}

Amazon Elastic Container Registry (ECR) jest prywatnym rejestrem obrazów kontenerów, zintegrowanym z~systemem
uprawnień IAM. W~zrealizowanym systemie utworzono po jednym repozytorium dla każdego mikroserwisu, z~włączonym
automatycznym skanowaniem obrazów pod kątem znanych podatności oraz polityką cyklu życia ograniczającą liczbę
przechowywanych wersji obrazu (koszt składowania rośnie z~każdą wersją, a~starsze obrazy szybko tracą użyteczność).
Węzły klastra pobierają obrazy na mocy roli IAM z~uprawnieniem tylko do odczytu, bez jakichkolwiek statycznych
poświadczeń.

Amazon Elastic Container Registry (ECR) jest prywatnym rejestrem obrazów kontenerów. W systemie utworzono repozytorium
dla serwisu ze skanowaniem obrazów pod kątem podatności. Ustawiono również cykl życia ograniczający liczbę przechowywanych wersji
obrazów. System ten zintegrowany jest z IAM. EKS za pomocą roli tylko do odczytu pobiera podczas deploymentu ostatnią wersje obrazu.

\subsection{Amazon Cognito - uwierzytelnianie użytkowników}
\label{subsec:cognito}

Amazon Cognito jest zarządzaną usługą tożsamości użytkowników końcowych. Jej centralnym pojęciem jest \emph{pula
użytkowników} (\emph{user pool}). Jest to katalog kont wraz z~mechanizmem rejestracji, logowania, polityki haseł
i~weryfikacji adresów e-mail. Po poprawnym uwierzytelnieniu wydaje podpisane cyfrowo tokeny JWT. Za ich pomocą
serwisy mogą obsłużyć logowanie. Każdy token może być zweryfikowany na podstawie publicznych kluczy udostępnianych
przez usługę. Serwis udostępnia również interfejs logowania (\emph{Hosted UI}). Dzięki niemu aplikacja webowa nie 
musi mieć formularza do logowania tylko korzysta z udostępnionego przez AWS. Szczegóły projektu uwierzytelniania 
przedstawiono w~rozdziale~\ref{chap:projekt-systemu}.

\subsection{Amazon S3 i CloudFront - hosting aplikacji webowej}
\label{subsec:s3-cloudfront}

Amazon Simple Storage Service (S3) jest obiektowym magazynem danych o~praktycznie nieograniczonej pojemności.
Amazon CloudFront jest siecią dystrybucji treści (CDN) serwującą zawartość z~lokalizacji brzegowych położonych
blisko użytkowników. Połączenie obu usług stanowi wzorcową architekturę hostingu aplikacji jednostronicowych
(SPA). Zbudowane statyczne artefakty przechowywane są w S3, a~CloudFront udostępnia je użytkownikom.
Bucket w którym przechowywana jest strona jest prywatny. Z tego powodu publiczny dostęp jest tylko z serwisu
CloudFront. Dostęp do niego CDN ma przez mechanizm emph{Origin Access Control} (OAC). Cloudfront
pełni w~systemie również role pojedynczego punktu wejścia dla API który kieruje ruch do load balancera.
Cloudfront również trzyma w pamięci podręcznej statyczne dane na lokalizacjach brzegowych, więc ogranicza on zapytania do S3
i przyspiesza dostarczenia strony. Konfigurację tę omówiono w~rozdziale~\ref{chap:infrastruktura}.

\subsection{Amazon VPC, ALB i architektura sieciowa}
\label{subsec:vpc-alb}

Amazon Virtual Private Cloud (VPC) udostępnia logicznie izolowaną sieć wirtualną o~adresacji definiowanej przez
użytkownika, dzieloną na podsieci rozmieszczone w~strefach dostępności. Podstawowym wzorcem zastosowanym
w zrealizowanym systemie jest podział na podsieci publiczne i prywatne. Sieć publiczna ma dostęp do internetu
dzięki temu może dostać ruch sieciowy. Sieć prywatna za to go nie ma i ruch do niej przechodzi przez bramę NAT.
Z tego powodu generalną praktyką jest wysyłanie ruchu internetowego do sieci publicznej i z sieci publicznej przez
NAT do prywatnej. Ruch sieciowy filtrowany jest przez \textbf{grupy bezpieczeństwa} (\emph{security groups}).
Są to stanowe zapory przypisywane zasobom, zezwalające wyłącznie na jawnie zdefiniowane przepływy.

Application Load Balancer (ALB) jest balancerem warstwy siódmej, trasującym żądania HTTP do zarejestrowanych celów
i~weryfikującym ich zdrowie. W~zrealizowanym systemie ALB przyjmuje ruch API od CloudFront i~kieruje go do węzłów
klastra. Dostęp do balancera ograniczono do zakresów adresowych CloudFront, a~żądania pozbawione ustawianego przez
CloudFront nagłówka weryfikacyjnego są odrzucane. Wymusza przepływ całego ruchu przez warstwę brzegową.

\subsection{IAM i Secrets Manager}
\label{subsec:iam-secrets-manager}

AWS Identity and Access Management (IAM) jest systemem zarządzania tożsamością i~uprawnieniami obejmującym całą
platformę. Uprawnienia działają na zasadzie polityk przypisywanych użytkownikom i i~\textbf{rolom}.
Role to tożsamości przyjmowane tymczasowo przez usługi i procesy. Działają na zasadzie najmniejszych 
uprawnień (\emph{principle of least privilege}). Każda tożsamość powinna posiadać wyłącznie uprawnienia niezbędne do swoich zadań. 
W~zrealizowanym systemie role IAM posiadają m.in.: płaszczyzna sterowania klastra, węzły robocze, sterownik EBS CSI.

AWS Secrets Manager uzupełnia IAM o~scentralizowane przechowywanie sekretów aplikacyjnych. Czyli haseł, kluczy
i~parametrów połączeń. Główną zaletą tego systemu jest szyfrowanie, kontrola dostępu i~audyt przechowanych sekretów.
Pozwala to nie trzymać haseł w kodzie źródłowym, a w pełni bezpiecznym środowisku.